\subsection{Endpoint, openness, and selector audit}
\label{subsec:new-endpoint-selector-audit}

The arguments above repeatedly pass between open role traces, closed maximal
reaches, and compact witness sets.  We record the precise endpoint statements
used in those passages.

\begin{lemma}[Open trace endpoint]
\label{lem:new-open-trace-endpoint}
Let $U$ be an open equilateral triangle and $T=\overline U$.  If
$T\cap[V,O]$ is the closed interval from $V$ to a point $P$, then
$P\notin U$.  If $Q$ is a noncentral point of a radial arm and $Q\in U$, then
$U$ has positive-length intersection with that arm.
\end{lemma}

\begin{proof}
The first assertion follows because $P$ is a boundary point of the closed
convex intersection.  For the second, a small Euclidean ball about $Q$ is
contained in $U$; its intersection with the arm is a relative open interval.
\end{proof}

\begin{lemma}[Compact set in an open unit triangle]
\label{lem:new-compact-open-shrink}
If a compact set $K$ is contained in an open unit equilateral triangle, then
\[
 \Lambda(K)<1.
\]
\end{lemma}

\begin{proof}
The distance from $K$ to each of the three side lines is positive.  Move the
three side lines inward by a common smaller distance.  Their intersection is
a closed equilateral triangle of side strictly below one and still contains
$K$.
\end{proof}

\begin{lemma}[Continuity of the enclosure number]
\label{lem:new-lambda-continuity}
For nonempty compact sets $K,L$,
\[
 |\Lambda(K)-\Lambda(L)|
 \le\frac3h d_H(K,L),
\]
where $d_H$ is the Hausdorff distance.
\end{lemma}

\begin{proof}
For every unit normal,
\[
 |h_K(n)-h_L(n)|\le d_H(K,L).
\]
Apply the support formula and take the minimum over orientations in both
directions.
\end{proof}

These lemmas justify all singleton-gap limits and every use of a radial trace
endpoint as a center-forced point.

\subsubsection{Branchwise proof of common-pair domination}

We expand the neighboring-capacity comparison in Lemma
\ref{lem:new-common-pair-domination}.  Assume first that $q\le p$.
The own-radial demand $1-q$ is admissible because the exact edge-normal
orientation has side
\[
 q+\max\{p,1-q\}=1.
\]
Thus
\[
 c_{\max}(p,q)\ge1-q.
 \tag{6.73a}
\]
For $C_+(p,q)$ there are three possibilities.

On the linear branch,
\[
 C_+(p,q)=1-q.
\]
On the plateau branch, $C_+=p(p)$, where the cubic root satisfies
\[
 f_p(1-q)
 =(1-q)^3-(p+2)(1-q)^2+2(p+1)(1-q)-1\ge0
\]
throughout the selected interval.  Since the cubic is increasing,
$p(p)\le1-q$.  On the radical branch, the capacity decreases from its
selected transition value and therefore is no larger than the plateau.  Hence
\[
 C_+(p,q)\le1-q.
\]
The reflected calculation gives the same bound for $C_-$ and treats the order
$p\le q$.  Combining with (6.73a) proves
\[
 C_\pm(p,q)\le c_{\max}(p,q).
\]

The same argument is valid for any actual pair $(a,b)$ with $a\ge p$ and
$b\ge q$, because the family of triangles containing the larger anchors is a
subfamily of the one containing the smaller anchors.  Thus the relevant
capacity is coordinatewise nonincreasing.

\subsubsection{Finite-caliper selector audit}

For a finite point set $K$, Proposition~\ref{prop:new-enclosure-gauge} reduces
$\Lambda(K)$ to finitely many orientations.  We spell out the reduction used
for every diagram and for the anisotropic one-gap set.  Let the convex hull
vertices be $P_0,\ldots,P_{r-1}$ in cyclic order.  On an angular interval on
which the three support points are fixed, write
\[
 H(\theta)=
 \sum_{j=0}^2
 \langle Q_j,n_{\theta+2\pi j/3}\rangle.
\]
Then
\[
 H''(\theta)=-H(\theta).
\]
The support sum is positive unless $K$ is a singleton at the origin.  Hence an
interior critical point is a strict maximum.  A minimum occurs when two
support points tie in one normal direction.  That normal is perpendicular to
a hull edge $P_iP_{i+1}$.  Therefore
\[
 \Lambda(K)=
 \min_{0\le i<r}
 \frac1h
 \sum_{j=0}^2
 h_K(\mathsf R^jn_i),
 \tag{6.73b}
\]
where $n_i$ is either unit normal to $P_iP_{i+1}$.

This finite-caliper identity has two consequences.  First, the enclosure
number is unchanged by replacing a witness set by its convex hull.  Second,
no numerical angular minimization is hidden in the proof: every terminal can
be checked at finitely many support contacts.  The complementary-gap theorem
uses this fact analytically; the zero-gap nine-point theorem uses it through
the support-arc certificate in the final appendix.

\subsubsection{Current gap-endpoint squeeze for the three-transverse witness}

Give a candidate CE1 or CE2 C triangle its signed variables and put
\[
 X=R-\delta,
 \qquad Q=\frac{\eta+\alpha+\delta}{2R}.
\]
The two actual gap endpoints lie strictly inside its open trace, so
\[
 B_0>2Q,
 \qquad A_1>X.
\]
The diameter inequality for the two incident boundary endpoints of $T_0$
then gives
\[
 A_0\le M_0(B_0)<M_0(2Q)
 =-Q+\sqrt{1-3Q^2}<1-Q.
\]
No point on $r_0$ occurs in this argument.  The current transverse witness
retains exactly $P_2,P_3,P_4$, and the complete calculation is
\zcref{lem:app-transverse-upper-squeeze}.

\subsubsection{Strictness at singleton gaps}

Suppose the V-gap on $e_{0,1}$ is the singleton $X(\ell)$.  Although its
length is zero, neither incident open V trace contains the point.  Boundary
locality therefore forces $X(\ell)\in U_C$.  Choose
$r_n>\ell$ decreasing to $\ell$, and reduce the incoming reach at $T_1$ from
$1-\ell$ to $1-r_n$.  Every local diameter inequality remains feasible, and
the strict supercritical pattern is unchanged for large $n$.  The non-singleton-gap
witness sets converge to the singleton witness set.  By Lemma
\ref{lem:new-lambda-continuity}, the inequality $\Lambda\ge1$ passes to the
limit.  Thus no proof in this section charges numerical gap length; every one
uses the actual missing point.

\subsubsection{Why a common disk is insufficient for $N_+=1$}

The one-gap $N_+=1$ terminal does not retain all six radial endpoints.  Its
exact witness is
\[
 K_{\mathrm{tr}}=
 \{O,M_0,X_0(B_0),X_0(1-A_1),P_2,P_3,P_4\}.
\]
The points $P_2,P_3,P_4$ supply, respectively, the transverse radial demands
used by the CE1 reverse return, while $P_2,P_4$ suffice for the CE2 threshold
alternative.  The upper squeeze uses only the two actual gap endpoints.  In
particular neither $P_0$ nor a nonexistent six-point radial catalogue enters
the canonical proof.
